\section{Transformer Attention: primero, definir con claridad qué se compara}

La clase no aplica la SDM directamente al Transformer; antes repasa la attention de principio a fin. El propósito es fijar qué se va a comparar: cómo la similitud entre query y key produce los pesos, qué hace softmax, por qué value es distinto de key y cómo se escribe en forma matricial el vector de actualización final.

\subsection{Semántica operativa de Q, K y V}

Esta subsección fija primero la semántica operativa de Q, K y V, de la que dependen las comparaciones geométricas posteriores. Después de la diapositiva de transición, el ejemplo de lenguaje trata la posición que se va a predecir como query, usa la key de cada token del contexto para el emparejamiento y trata la value como el contenido que realmente se transfiere a la salida. La diapositiva de Q/K/V divide el proceso en generación de proyecciones, comparación query--key, normalización y ponderación de las values; la diapositiva siguiente muestra cómo los pesos normalizados multiplican cada value y luego se suman.

\lecturefigure{slide-13-transformer-attention-overview.jpg}{Transición a la sección de Transformer attention}{Diapositiva V013 recuperada de la grabación oficial; 00:07:20}

\lecturefigure{slide-14-attention-example.jpg}{Ejemplo de attention en un modelo de lenguaje}{Diapositiva V014 recuperada de la grabación oficial; 00:08:02}

\lecturefigure{slide-15-qkv-steps.jpg}{Los cuatro pasos de cálculo con keys, queries y values}{Diapositiva V015 recuperada de la grabación oficial; 00:08:40}

\lecturefigure{slide-16-weighted-value-sum.jpg}{Cómo los pesos de softmax agregan los value vectors}{Diapositiva V016 recuperada de la grabación oficial; 00:09:00}

\begin{importantbox}{Key y value no deben confundirse}
La key responde «qué tan relacionado está este recuerdo con la query», y la value responde «si está relacionado, qué debe recuperarse». Al separarlas, el modelo puede emparejar por sintaxis o por identidad de entidad y, aun así, devolver características distintas que sean útiles para el siguiente token; esto también corresponde a la hetero-association de la SDM, en la que address y stored content pueden separarse.
\end{importantbox}

\subsection{Ejemplo ficticio y ecuación de actualización completa}

La diapositiva Fictional example usa tokens como ``cat'' y ``sat'' para mostrar que una key muy relacionada puede devolver una value que contiene información semántica y sintáctica. A continuación, la diapositiva Full query update escribe juntas la proyección y la actualización. Sea $X$ la matriz de tokens de entrada, $x_q$ el vector de la posición de consulta y $W_Q,W_K,W_V$ las matrices de proyección; entonces
$$
q=W_Qx_q,\quad K=XW_K,\quad V=XW_V,
$$
$$
\operatorname{Attn}(q,K,V)
=V^{\top}\operatorname{softmax}\!\left(\beta Kq\right).
$$
Aquí $q$ es la query, cada fila de $K$ es una key, cada fila de $V$ es una value y $\beta$ es el coeficiente de temperatura o escala. Si se usa scaled dot-product, lo habitual es $\beta=1/\sqrt{d_k}$.

\lecturefigure{slide-17-fictional-example.jpg}{Ejemplo ficticio: al emparejar una key se devuelven características de value distintas}{Diapositiva V017 recuperada de la grabación oficial; 00:10:14}

\lecturefigure{slide-18-full-query-update.jpg}{Expresión matricial de la query update completa}{Diapositiva V018 recuperada de la grabación oficial; 00:10:30}

\subsection{Dot product, softmax y value summation}

Tres diapositivas amplían, cada una, una operación central de la attention. La diapositiva de dot product explica que la similitud proviene de $k_i^{\top}q$; la de softmax explica que la exponenciación amplía las diferencias entre puntuaciones altas y normaliza; la de summation explica que la salida es una combinación convexa de las values. Lo esencial de softmax no es la etiqueta de «producir probabilidades», sino que los pesos varían exponencialmente con la similitud, que es justamente la forma que la sección siguiente busca recuperar a partir de la intersección de esferas en alta dimensión.

\lecturefigure{slide-19-query-key-dot-product.jpg}{El dot product query--key forma los similarity logits}{Diapositiva V019 recuperada de la grabación oficial; 00:10:50}

\lecturefigure{slide-20-softmax.jpg}{Softmax amplifica los términos de alta similitud y completa la normalización}{Diapositiva V020 recuperada de la grabación oficial; 00:11:44}

\lecturefigure{slide-21-summation-over-values.jpg}{Suma de los value vectors según los pesos de atención}{Diapositiva V021 recuperada de la grabación oficial; 00:12:02}

\lecturefigure{slide-22-full-attention-equation.jpg}{Flujo completo de la attention y ecuación de actualización final}{Diapositiva V022 recuperada de la grabación oficial; 00:12:18}

Para la $i$-ésima key, el peso de softmax es
$$
\alpha_i(q)=\frac{\exp(\beta k_i^{\top}q)}
{\sum_j\exp(\beta k_j^{\top}q)},
\qquad
y(q)=\sum_i\alpha_i(q)v_i.
$$
Aquí $\alpha_i$ es el normalized retrieval weight y $v_i$ es el contenido del $i$-ésimo recuerdo. Al compararla con la SDM, la pregunta se convierte en: si el sphere-intersection count puede escribirse como $\exp(\text{similarity})$ y si, tras normalizar, se obtiene el mismo $\alpha_i$.

\begin{lstlisting}[caption={Scaled dot-product attention (ilustración con una sola query)}]
def attention(query, keys, values, beta):
    logits = beta * (keys @ query)
    weights = softmax(logits)
    return weights @ values
\end{lstlisting}

\teachervoice{El ponente advierte en particular que «what you pay attention to» y «what you get out» pueden ser distintos: la query y las keys determinan la relevancia, y las values determinan el contenido recuperado. Esta idea es esencial para separar más adelante la cerebellar activation pathway de la write pathway, y también se conecta con la content/structure factorization de la segunda clase.}

\subsection{Resumen de la sección}

La attention es una memoria direccionable por contenido: la query y las keys producen similitudes, softmax las convierte en pesos normalizados exponencialmente y las values se recuperan ponderadas. Siempre que la sphere intersection de la SDM decaiga exponencialmente con la distancia y la distancia pueda reescribirse como producto punto, ambas operaciones de lectura comparten la misma estructura de cálculo.
